% Add "answers" option to print the answers
\documentclass[addpoints]{exam}

\usepackage[utf8]{inputenc}
\usepackage{calc}
\usepackage[dvipsnames]{xcolor}
\usepackage{ifthen}
\usepackage{tcolorbox}

% \makebox[\textwidth]{Name:\enspace\hrulefill}
%
% \vspace{5mm}
%\makebox[\textwidth]{Instructor’s name:\enspace\hrulefill}

\definecolor{MyLightGray}{gray}{0.60}

\newcommand{\truefalse}
{		\begin{checkboxes}
			 \choice True
			 \choice False. Explanation:
		\end{checkboxes}
		\vspace{10pt}
}


% Check if 'answers' is an option
\newcommand{\ifanswers}[2]{%
  \ifthenelse{\boolean{exam@answers}}{#1}{#2}%
}



\newcommand{\answerspace}[1]
{
	\ifprintanswers
	\else
		\setlength{\mylength}{10pt}%
		% \vspace{#1\mylength}
		% Make the grey box
		\par\nobreak\noindent
		\makebox[\linewidth][l]{%
			% set inner padding to zero, so only the frame shows
			{\setlength{\fboxsep}{0pt}%
			% draw a color frame, white (or transparent) inside
			\fcolorbox{MyLightGray}{white}{%
			% the inner box has zero width content but fixed height
			\parbox[t][#1\mylength][c]{%
				% subtract twice the frame‐thickness from the full width
				\dimexpr\linewidth-2\fboxrule\relax
			}{\strut}%
			}%
			}%
		}%
		\par
		\vspace{3\mylength}

  	\fi
}
\newlength{\mylength}

\newcommand{\code}[1]{{\texttt{#1}}}


%\usepackage{listings}
\usepackage{color}
\usepackage{lastpage}
\usepackage{verbatim}
\usepackage{fancyvrb}
\usepackage{graphicx}
\usepackage{float}
\usepackage{minted}
\usepackage{multicol}

\newcommand{\includecode}[2][c]{\lstinputlisting[caption=, escapechar=, language=#1]{#2}}
\RecustomVerbatimCommand{\VerbatimInput}{VerbatimInput}{}

\renewcommand{\solutiontitle}{\noindent\textbf{Solution sketch: }}

\definecolor{SolutionColor}{rgb}{0.8,0.9,1}

\cfoot{Page \thepage\ of \pageref{LastPage}}

\rhead{Initials: \rule{20mm}{0.4pt}}



%%%%%%%%%%%%%%%%%%%%%%%%%%%%%%%%%%%%%%%%
\begin{document}
\fvset{numbers=left,numbersep=3pt}
\thispagestyle{empty}

\begin{center}
	{\LARGE CS 408 Fall 2026 -- Mutation Testing Quiz \\}
\fbox{\fbox{\parbox{5.5in}{\centering
	\begin{itemize}
        \item Try to come up with a solution on your own first, but
          after that talk to your friends and neighbors.

        \item Full points will be rewarded for \textbf{any} (reasonable)
          attempt.
        \end{itemize}
	}}}
\end{center}

\vspace{5mm}
\makebox[\textwidth]{\textbf{Name (please print)}: \rule{100mm}{0.8pt} \hfill}

\vspace{10mm}
\makebox[\textwidth]{\textbf{PUID}: \rule{100mm}{0.8pt} \hfill}


\begin{questions}

\question The method below charges \texttt{days} $\times$
  \texttt{rate}, less a \$10 discount for rentals longer than a week.
  The fee is never negative.

\begin{minted}{java}
static int rentalFee(int days, int rate) {
  int total = days * rate;
  if (days > 7) {
    total = total - 10;
  }
  if (total < 0) {
    total = 0;
  }
  return total;
}
\end{minted}

  Here are three mutation operators:

  \begin{description}
  \item[Conditionals Boundary] replace a relational operator with its
    boundary variant: \texttt{>} $\leftrightarrow$ \texttt{>=},
    \texttt{<} $\leftrightarrow$ \texttt{<=}
  \item[Math] replace one arithmetic operator with another:
    \texttt{*} $\rightarrow$ \texttt{+}, \texttt{-} $\rightarrow$
    \texttt{+}, and so on
  \item[Inline Constant] replace a numeric literal with a different one
  \end{description}

  Using any one of the three operators, write down a first-order
    mutant\footnote{ A \textbf{first-order} mutant applies exactly one
      such change.} of \mintinline{java}{rentalFee()}. Say which
    operator you used and which line you changed. Then give a test
    case that \textbf{kills} it, along with both results.
    \answerspace{14} %
    \begin{solution} %
      Many correct answers. For instance, \textbf{Math} on line 2,
      changing \mintinline{java}{days * rate} to
      \mintinline{java}{days + rate}:
      \mintinline{java}{rentalFee(3, 4)} returns 12 in the original and
      7 in the mutant.

      Other killable mutants:

      \begin{tabular}{lll}
        \textbf{operator} & \textbf{change} & \textbf{killed by} \\
        \hline
        Boundary & line 3: \texttt{days > 7} $\to$ \texttt{days >= 7}
          & \texttt{(7, 10)}: 70 vs.\ 60 \\
        Math & line 2: \texttt{days * rate} $\to$ \texttt{days + rate}
          & \texttt{(3, 4)}: 12 vs.\ 7 \\
        Math & line 4: \texttt{total - 10} $\to$ \texttt{total + 10}
          & \texttt{(8, 10)}: 70 vs.\ 90 \\
        Constant & line 3: \texttt{7} $\to$ \texttt{1}
          & \texttt{(2, 10)}: 20 vs.\ 10 \\
        Constant & line 4: \texttt{10} $\to$ \texttt{1}
          & \texttt{(8, 10)}: 70 vs.\ 79 \\
      \end{tabular}
    \end{solution}

  % \part Now write down a \emph{different} first-order mutant, using a
  %   \textbf{different} operator than in part (a), and give a test case
  %   that kills it.
  %   \answerspace{14} %
  %   \begin{solution} %
  %     Any other row of the table above. The point is that each operator
  %     is a rule that can be applied at several places in the program ---
  %     a real tool applies all of them everywhere, which is where the
  %     thousands of mutants come from.
  %   \end{solution}

  % \part There is one mutant you can build with the \textbf{Conditionals
  %   Boundary} operator that \emph{no} test case can kill. Which is it,
  %   and why?
  %   \answerspace{12} %
  %   \begin{solution} %
  %     Line 6: \mintinline{java}{total < 0} $\to$
  %     \mintinline{java}{total <= 0}. The two versions behave
  %     differently only when \texttt{total} is exactly 0, and in that
  %     case the mutant executes \mintinline{java}{total = 0;}, which
  %     changes nothing. The mutant is \textbf{equivalent} to the
  %     original.

  %     We cannot detect this automatically --- deciding whether two
  %     programs are equivalent is undecidable --- so a kill rate of
  %     100\% is not achievable here, and surviving mutants have to be
  %     triaged by hand.
  %   \end{solution}

  % \end{parts}

\end{questions}

\end{document}
